%% notes-tex/figure-3-gate/sections/2-mockups.tex
%% SOURCE: figures are exported from notes/assets/drawio/fig3-mockup-*.drawio.svg, written by
%% experiments/019-simb-multimodal/scripts/fig3_mockup_drawio.py, which embeds the panels of
%% experiments/019-simb-multimodal/scripts/fig3_mockup_panels.py. Caption numbers are read
%% off the same result files; none is typed into a figure. The source tables under the
%% figures are hand-authored lists of the files and keys each panel function reads.

\section{The mockups\secstatus{todo}}
\label{sec:mockups}

Result files are under \file{experiments/019-simb-multimodal/results/} unless a path is
given. Each caption gives, per panel, what is plotted, the statistic and its scoring rule,
the sample size and the status; the table under each figure gives the files and keys.

\clearpage
\tcfig{figures/fig3-mockup-A-conditioning.pdf}{\textbf{Ridge from one measured modality
reaches a held-out Pearson of 0.17 to 0.23 on all features of another, and the conditioned
model (d) has trained 160 of 1,200 epochs.}
\textbf{a}, The three single-deletion panels and their pairwise overlaps. Counts of distinct
deletion strains; no statistic. Ready.
\textbf{b}, Replicate ceiling (pale; dotted, the HIS3-replicate proteome estimator), best
genotype-only model (solid; percent of ceiling), best linear baseline (diagonal hatch) and
best nearest-neighbor baseline (cross hatch), as Pearson per feature across held-out
strains. Model: expression and morphology are the rolling maximum of one best run over 1,353
and 183 runs, the proteome is the mean over 4 split seeds. Baselines: mean validation score
over 4 split seeds of the best key (morphology, 1 split seed); no linear baseline exists for
morphology. Every non-ceiling bar is a maximum over runs or keys selected on the score it
reports, so it is biased upward, and the bars sit on different splits. Partial.
\textbf{c}, Out-of-fold ridge between measured modalities on the deletions both panels
measured (1,349, 4,314 and 1,438 strains; 5 folds): median held-out Pearson per feature;
\emph{moving} is the 116 morphology features with replicate reliability of at least 0.5. The
proteome and expression side is a mean per feature from a review file. Partial.
\textbf{d}, Round v20, conditioned arm minus unconditioned control: difference of the mean
validation Pearson over the latest 20 epochs of each conditioned run and the same epochs of
its control; 2 of 12 split seeds, 8 runs, epochs 157 to 163 of 1,200. Not a result. Partial.
\textbf{e}, Conditioned morphology. Not run. Not done.
\textbf{f}, Linear and nearest-neighbor baselines and the v19 single-head model per split
seed (expression 11, proteome 4). Baselines: one validation Pearson per split seed, best
embedding selected on that score (biased upward). Model: mean validation Pearson over the
registered epoch window (unselected). A split seed is one train, validation and test
partition with one initialization seed; baselines and model share the seed number and not
the store, so their validation strains differ. Partial.
\textbf{g}, Per graph, the AUC that two deletions adjacent on the graph give more similar
phenotypes than a background pair (filled), against a degree-preserving rewired graph
(open); 60,000 pairs per point, 1,482 expression and 4,718 morphology strains; no proteome
strand is on file. Partial.
\textbf{h}, Spearman between two panels' gene-by-gene correlation matrices, with the number
of shared genes; Messner is the proteome, the others mRNA. One value per pair, no
selection. Ready.}{fig:mockup-a}

\clearpage
{\footnotesize\noindent
\begin{tabular}{@{}p{8mm}p{170mm}@{}}
\toprule
\multicolumn{2}{@{}l}{\textbf{Sources for \cref{fig:mockup-a}}} \\
\midrule
a & \file{fig3_overlap_census.json} (\texttt{kemmeren\_unique}, \texttt{ohya\_unique}); \file{fig3_proteome_build_census.json}; \texttt{n\_shared\_strains} of the two covariation files \\
b & \file{expression_ceiling_replicate.json}, \file{proteome_ceiling_replicate.json}, \file{morphology_noise_ceiling.json}, \file{round_leaderboards.csv}; \file{baselines_embedding_study.json} (ridge $0.113$, $0.066$); \file{graph_retrieval_baseline.json} (nearest neighbor $0.220$, $0.112$, STRING experimental graph key); \file{knn_embedding_probe.json} ($0.100$) \\
c & \file{proteome_morphology_covariation.json} (and its \texttt{context} block), \file{expression_morphology_covariation.json} \\
d & \file{joint_checkpoint_readout.json}, \texttt{v20\_partial.runs} \\
f & \file{expression_baselines_split/}, \file{baselines_split_fig3_proteome/}; \file{joint_checkpoint_readout.json}, \texttt{v19.runs} \\
g & \file{graph_prior_probe.json}, \texttt{pair\_form}, one step \\
h & \file{experiments/028-knockout-expression/results/proteome_expression_covariation.json}, \texttt{gene\_covariation} \\
\bottomrule
\end{tabular}}

\clearpage
\tcfig{figures/fig3-mockup-B-genotype-only.pdf}{\textbf{From the deletion alone the model
reaches 13 to 29\,\% of the replicate ceiling, and the best nearest-neighbor baseline is
within 0.01 of it on expression and above it on the proteome and morphology.}
\textbf{a}, \textbf{b}, As \cref{fig:mockup-a}a and b. Ready; partial.
\textbf{c}, Round v19, joint arm minus single-head arm per split seed: difference of the
mean validation Pearson per feature over the registered window (expression epochs 1,000 to
1,200, proteome 200 to 400), an unselected statistic. 11 of 12 split seeds, each a different
train, validation and test partition with one initialization seed; 1,099 to 1,108 training
and 119 to 126 validation strains on the four seeds on file. Expression: mean $-0.001$, 5 of
11 above zero, one-sided $t$ $p = 0.55$. Proteome: mean $-0.017$, 1 of 11, $p = 0.64$
against the margin $-0.015$ (dashed). Neither registered test rejects. Ready; split seed 11
was unfinished.
\textbf{d}, As \cref{fig:mockup-a}f at half width, same caveat. Partial.
\textbf{e}, As \cref{fig:mockup-a}g at wide width. Partial.
\textbf{f}, How much the gene representation matters to the two baselines: validation
Pearson per feature, mean and s.d. over 4 split seeds, of the ridge baseline (squares) and
the nearest-neighbor baseline (triangles), with rank and neighbor count selected on
validation inside each seed. Five of 28 representations by rule, each the top of its class
on the mean of the four validation means: best single language model (ProtT5), the model's
input composite (bold), best other composite, codon frequency, best random control. The
bold row is the study's \texttt{model\_stack}, the concatenation
\texttt{fudt\_upstream}, \texttt{calm}, \texttt{prot\_T5\_all}, \texttt{fudt\_downstream}
that the v19 configuration inherits as \texttt{node\_embeddings}. Ready.}{fig:mockup-b}

\clearpage
{\footnotesize\noindent
\begin{tabular}{@{}p{8mm}p{170mm}@{}}
\toprule
\multicolumn{2}{@{}l}{\textbf{Sources for \cref{fig:mockup-b}}} \\
\midrule
c & \file{joint_checkpoint_readout.json}, \texttt{v19.tests} (\texttt{H1a\_expression\_superiority}, \texttt{H1b\_proteome\_noninferiority}); strain counts from \file{split_gene_overlap_audit.json} \\
f & \file{baselines_embedding_study.json}; the composite is defined in \file{experiments/019-simb-multimodal/scripts/expression_baselines_split.py} and matched to \file{conf/cgt_expr_v13_split.yaml} \\
\bottomrule
\end{tabular}}

\clearpage
\tcfig{figures/fig3-mockup-SI-1.pdf}{\textbf{In round v19 the validation loss reaches its
minimum before epoch 200 on every run, and the proteome score peaks later on every run.}
\textbf{a}, How a split seed becomes the strains of a run, drawn from the code. The store's
4,681 records (one record is one deleted gene set) are shuffled within each phenotype label
and each perturbation count and cut 80, 10, 10 under the split seed; each side then keeps
the records carrying both labels, 1,349 in all, so the both-label fractions are near and not
exactly 80, 10, 10. Counts for split seeds 0 to 3. No statistic. Partial: seeds 4 to 11 are
not on file.
\textbf{b}, Proteome head, split seed 5, one run per arm: the 5-epoch centered mean of the
validation loss (top) and the validation Pearson per feature (bottom; pale raw, dark its
5-epoch mean) against epoch. Triangles, loss minimum; circles, maximum of the 5-epoch mean,
an upward-biased order statistic; shaded, the registered window. The joint loss sums two
heads, so the two loss levels are not comparable. Split seed 5 is chosen by rule: its
joint-minus-single proteome difference ($-0.012$) is the median over the 11 complete split
seeds. About 1,100 training and 125 validation strains. Ready.
\textbf{c}, The same for the expression head of that split seed, window epochs 1,000 to
1,200. Ready.
\textbf{d}, Proteome head, all 11 complete split seeds, 22 runs: epoch of the loss minimum
(triangles) and of the score peak (circles). Ready.
\textbf{e}, The window scores behind \cref{fig:mockup-b}c: mean validation Pearson over the
registered window, single-head and joint arm, 11 split seeds, 33 runs. Unselected.
Ready.}{fig:mockup-si-1}

\clearpage
{\footnotesize\noindent
\begin{tabular}{@{}p{8mm}p{170mm}@{}}
\toprule
\multicolumn{2}{@{}l}{\textbf{Sources for \cref{fig:mockup-si-1}}} \\
\midrule
a & \file{torchcell/datamodules/cell.py} (\texttt{\_compute\_and\_save\_index}); \file{experiments/019-simb-multimodal/scripts/train_cgt_multitask.py} (\texttt{require\_modalities}); counts from \file{split_gene_overlap_audit.json} \\
b, c & \file{joint_checkpoint_curves.csv} (41,129 rows: run, arm, split, epoch, loss, its 5-epoch mean, each head's Pearson), written by \file{experiments/019-simb-multimodal/scripts/joint_checkpoint_readout.py}; marker epochs from \file{joint_checkpoint_readout.json} \\
d, e & \file{joint_checkpoint_readout.json}, \texttt{v19.runs} \\
\bottomrule
\end{tabular}}

\clearpage
\tcfig{figures/fig3-mockup-SI-2.pdf}{\textbf{The model's input composite scores below
ProtT5 alone in both baselines on both modalities, and the best score of a run does not
track its parameter count.}
\textbf{a}, All 28 gene representations of the baseline study, by family: validation
Pearson per feature, mean and s.d. over 4 split seeds, ridge (squares) and nearest neighbor
(triangles), expression (1,244 to 1,253 training strains) and proteome (3,581); rank and neighbor
count selected on validation inside each seed. Bold, the model's input composite. Ready.
\textbf{b}, The six directed pairs of \cref{fig:mockup-a}c as bars with the strain-permuted
null (black line, on file for four pairs) and the shared-strain count. Partial.
\textbf{c}, The graph as a retrieval key: validation Pearson of a neighbor average keyed on
each graph's adjacency row minus the same with the rewired graph, mean and s.d. over 4 split
seeds, expression and proteome. Neighbor count selected on validation. Ready.
\textbf{d}, Parameter count against the rolling maximum of the validation Pearson, one point
per expression-strand run (1,182 runs with both values; open, collapsed runs). The rolling
maximum is biased upward and the runs differ in epochs run. The Pearson between
$\log_{10}$ parameters and score is $-0.001$. Dashed, the v19 model at 6.67 million
parameters for about 1,100 training strains, near 6,000 parameters per strain; its score is
a window mean and is not placed on this axis. Partial.
\textbf{e}, The perturbation operator in the notation of manuscript Fig.~1f; no data. Left,
as drawn there. Right, as the code computes it and as the Methods write it: a softmax over
the perturbed genes, which with one perturbed gene gives weight 1 to every gene. Partial,
an author decision.}{fig:mockup-si-2}

\clearpage
{\footnotesize\noindent
\begin{tabular}{@{}p{8mm}p{170mm}@{}}
\toprule
\multicolumn{2}{@{}l}{\textbf{Sources for \cref{fig:mockup-si-2}}} \\
\midrule
a & \file{baselines_embedding_study.json} \\
b & as \cref{fig:mockup-a}c, plus \texttt{triangle} of \file{joint_checkpoint_readout.json} for one strain count \\
c & \file{graph_retrieval_baseline.json}, rounds \texttt{v13} and \texttt{v14}, \texttt{summary.paired} \\
d & \file{round_leaderboards.csv} (\texttt{total\_param\_count}, \texttt{primary\_roll\_max}); \texttt{total\_param\_count} of \texttt{v19.runs} in \file{joint_checkpoint_readout.json}; \file{split_gene_overlap_audit.json} \\
e & \file{torchcell/models/equivariant_cell_graph_transformer.py} (\texttt{EquivariantPerturbationTransform}); \file{paper/nature-biotech/sections/methods.tex} \\
\bottomrule
\end{tabular}}
